\documentclass[11pt]{article}
\usepackage[margin=1in]{geometry}
\usepackage{amsmath,amssymb,amsthm,hyperref}
\newtheorem{definition}{Definition}
\newtheorem{theorem}{Theorem}
\newtheorem{lemma}{Lemma}
\newtheorem{corollary}{Corollary}
\newtheorem{obligation}{Open obligation}
\newcommand{\SAT}{\mathsf{SAT}}
\newcommand{\PP}{\mathsf P}
\newcommand{\NP}{\mathsf{NP}}
\title{LOWER\_BOUND-0001: Universal Quantifiers, Bounded Lists, and the Oracle Barrier}
\author{P versus NP project}
\date{10 October 2026}
\begin{document}
\maketitle
\paragraph{Status.}
This manuscript proves elementary model-audit lemmas and a matched-oracle
extension of the prefix-search immunity argument in Hirahara
\cite{hirahara2020}. It closes a barrier obligation deferred in project
PVG-4-0001, not the existence obligation for PFCERT. No separation of
$\PP$ and $\NP$ is claimed. These are human-readable formal proofs,
not proof-assistant verification.

\section{What history distinguishability actually proves}
\begin{definition}[Summary-only continuation]
Fix a set $H$ of histories, a common query set $Q$, and a response function
$f:H\times Q\to\{0,1\}$. A summary-only implementation consists of
$s:H\to\{0,1\}^{b}$ and one function
$E:\{0,1\}^{b}\times Q\to\{0,1\}$ satisfying
$E(s(h),q)=f(h,q)$. No history-dependent source, advice, oracle, or
unrecorded state is available to $E$.
\end{definition}
\begin{lemma}[Summary bound]
If $H_0\subseteq H$ has $M$ pairwise distinct response rows
$f(h,\cdot)$, then every summary-only implementation has
$b\ge\lceil\log_2M\rceil$.
\end{lemma}
\begin{proof}
If $s(h)=s(h')$, the common evaluator has the same answer for every
query, so the two rows are equal. Thus $s$ is injective on $H_0$.
\end{proof}
This proves neither $M$ time steps nor a bound for a model allowing
source rereading. For example, let $H=\{0,1\}^{k}$ and
$f(x,i)=x_i$. All $2^k$ rows differ, but if $x$ remains on the read-only
random-access input tape, each query requires just its indexed input
read and ordinary address handling. There need be no retained $k$-bit
copy of $x$ in work space.

\begin{lemma}[Endpoint bypass, recovered]
Assume $A$ is an exact polynomial-time SAT decider. If the sole required
output on a formula $F$ is $\SAT(F)$ and the ordinary model permits
access to $F$ and execution of $A$, then there is a correct implementation
that ignores every wrapper-created history and returns $A(F)$.
\end{lemma}
\begin{proof}
The stated implementation has the required output and the assumed
polynomial time bound by the definition of $A$.
\end{proof}
Consequently, an auxiliary witness being unique, or a false history
destroying satisfying extensions of a wrapper CNF, does not prove that
an arbitrary decision algorithm constructs that witness or history.

\section{The current all-input gap}
Let $\mathcal A$ be the set of simultaneous SAT time-space exponent
pairs for one uniform all-length deterministic machine on the declared
encoding and random-access model. The project has proved a conditional
recursive composition rule
\[
 C=\max\left\{1,r,\frac{v-h}{1-h}\right\},\qquad
 P=\max\{p,h\},
\]
for a uniformly hereditary, fully charged all-input phase with local
time $N^{r+o(1)}$, aggregate executed child volume $N^{v+o(1)}$,
maximum child length $N^{h+o(1)}$, $h<1$, and local state
$N^{p+o(1)}$. Repeated children are counted with multiplicity.
The proof and full contract are in FORMAL-SPINE1 and
RECOVERY-STRENGTHEN7, identified in the companion audit report.

\begin{lemma}[Existential phase equivalence, recovered]
Suppose a phase contract permits zero children and local direct decision
at resource bounds $(C,P)$, and recursive composition converts every
such phase to a SAT algorithm with the same bounds. Then existence of
an all-input phase is equivalent to $(C,P)\in\mathcal A$.
\end{lemma}
\begin{proof}
The forward implication is the assumed composition theorem.
For the reverse implication, execute the algorithm locally and issue no
children. The child restrictions hold vacuously.
\end{proof}
A positive child count alone does not repair this: one can issue and
ignore a constant SAT instance.

\begin{obligation}[Source-visible all-input construction]
Prove, under the assumption that SAT has a polynomial-time decider, the
existence of a uniform all-length implementation satisfying a favorable
fully charged hereditary phase contract on every root and actual
unresolved descendant. A sufficient stronger target is an explicit
compiler from a clocked SAT decider to a PRE-A construction that builds
strictly smaller, semantically essential children before any parent-scale
semantic call. Acquisition failures, replay, input access, current child
buffers, and all ancestor state must be charged in the same machine.
\end{obligation}
This obligation is separation-strength for contracts whose endpoint is
already forbidden. The following verified endpoint is a checkpoint,
not a solution to the obligation.

Fortnow--Lipton--van Melkebeek--Viglas \cite{flmv2004}, Corollary 5.8
and the local SAT reduction, exclude $(t,s)$ with strict margins
\[
 t<c,\qquad
 s<d(c):=\frac12\left(\frac{c+2}{c^2}-c\right).
\]
For $c=1201/1000$,
\[
 d(c)=\frac12+\frac{26275399}{2884802000},
\]
so $(6/5,1/2)$ is forbidden. Matching representation, local reduction,
and simultaneous resource bounds are essential. This is not an
unrestricted polynomial-time SAT lower bound.

\section{Matched oracle and machine conventions}
Fix an arbitrary oracle $O\subseteq\{0,1\}^{*}$ and a universal binary-program
oracle machine $U^O$. Every computation and verifier in this section uses
the same $O$. A program has at most one final output, produced when it
halts. Query writing and ordinary tape work are counted; the oracle
answer step has the usual unit cost.

We use a plain universal machine and require efficient universality:
for each fixed plain-description oracle machine $M$
there is a compiler prefix of constant length independent of its input,
and simulation by $U^O$ has polynomial overhead in input length and
running time, uniformly in $O$. No constant-overhead claim for simulating
arbitrary plain-description machines by a prefix-free universal machine
is made. A self-delimiting encoding of
$n$ of length at most $2\lceil\log_2n\rceil+O(1)$ is used.

Let $t:\mathbb N\to\mathbb N$ eventually dominate every polynomial.
Define
\[
 K_U^{t,O}(x)=\min\{|p|:U^O(p)\text{ halts with final output }x
                   \text{ within }t(|x|)\text{ steps}\},
\]
with minimum $+\infty$ when no such program exists.

\begin{theorem}[Relativized immunity with logarithmic deficiency]
Let $b(n)$ be a nonnegative integer with
$b(n)\le c\log_2(n+2)$ for some constant $c$. Then
\[
 R_b^{t,O}=\{x:K_U^{t,O}(x)\ge |x|-b(|x|)\}
\]
is $\PP^O$-immune.
\end{theorem}
\begin{proof}
Counting programs of length strictly less than $n-b(n)$ shows that,
for every sufficiently large $n$, at least one $n$-bit string lies in
$R_b^{t,O}$. Thus the set is infinite.

Suppose an infinite $L\in\PP^O$ were contained in $R_b^{t,O}$.
Fix its deterministic polynomial-time decider $M^O$.
Choose a positive integer constant $a>c+4$ and set
$m(n)=\lceil a\log_2(n+2)\rceil$.
All sufficiently large $n$ have $m(n)<n$; ignore the finitely many others.

For any sufficiently large length $n$ at which $L$ is nonempty, choose
one $y\in L\cap\{0,1\}^n$ and let $u$ be its first $n-m(n)$ bits.
A fixed oracle program, given encoded $n$ and $u$, enumerates all
$m(n)$-bit suffixes $v$, tests $M^O(uv)$, and prints the first accepted
$uv$. Such a suffix exists. Its running time, including universal
simulation, is polynomial in $n$: there are
$2^{m(n)}=n^{O(1)}$ tests of a fixed polynomial-time machine. This time
is eventually at most $t(n)$.

The program description has length at most
\[
 n-m(n)+2\lceil\log_2 n\rceil+O(1)
 <n-b(n)
\]
for all sufficiently large $n$. The selected output $x=uv$ belongs
to $L$, but has $K_U^{t,O}(x)<n-b(n)$, a contradiction.
\end{proof}

The proof is a direct matched-oracle adaptation and logarithmic-deficiency
extension of the published prefix-search argument, not a new
nonrelativizing mechanism.

\section{A generalized bounded-list target}
Let $L(n)\ge1$ be an integer, polynomially bounded and polynomial-time
computable. Define
\[
 \ell(n)=\lceil\log_2 L(n)\rceil,\qquad
 k(n)=n-\ell(n)-1.
\]
Only sufficiently large $n$, for which $k(n)\ge0$, are relevant.

\begin{definition}[Polynomially witnessed endpoint list cover]
One fixed deterministic polynomial-time oracle verifier
$V^O(1^n,p,x,\pi)$ and one fixed polynomial $q$ satisfy:
\begin{description}
\item[] \emph{Witness convention.} Every accepted witness is restricted
to $|\pi|\le q(n)$; outputs have $|x|=n$, and $|p|\le k(n)$.
\item[] \emph{Completeness.} If $U^O(p)$ halts with final output
$x\in\{0,1\}^n$ within $t(n)$ steps, then some permitted $\pi$
satisfies $V^O(1^n,p,x,\pi)=1$.
\item[] \emph{Global list soundness.} For each fixed $(n,p)$,
\[
 \left|\{x\in\{0,1\}^n:\exists \pi,\ |\pi|\le q(n),\
                 V^O(1^n,p,x,\pi)=1\}\right|\le L(n).
\]
This bounds the union over all permitted witnesses, not the output of
one branch or one run.
\end{description}
\end{definition}

\begin{theorem}[List-cover conditional separation]
If this endpoint list cover exists, then $\PP^O\ne\NP^O$.
\end{theorem}
\begin{proof}
Let $E$ be the language whose $n$-bit members are the accepted outputs
over all $p$ of length at most $k(n)$ and permitted witnesses.
Because $k(n)\le n$ and $q$ is polynomial, $E\in\NP^O$.
There are $2^{k(n)+1}-1$ possible programs, so
\[
 |E\cap\{0,1\}^n|
 \le L(n)(2^{k(n)+1}-1)
 \le 2^n-L(n)<2^n.
\]
Thus $\overline E$ has a member at every sufficiently large length.
By completeness, any such member satisfies
\[
 K_U^{t,O}(x)>k(n),\qquad
 K_U^{t,O}(x)\ge n-\ell(n).
\]
After deleting finitely many short lengths, $\overline E$ is an infinite
subset of $R_\ell^{t,O}$. Since $L$ is polynomially bounded,
$\ell(n)=O(\log n)$, so the preceding immunity theorem applies.
If $\PP^O=\NP^O$, then $E$ and its complement are in $\PP^O$,
contradicting that immunity.
\end{proof}

\begin{corollary}[Historical two-valued case]
Taking $L(n)=2$ gives $\ell(n)=1$ and $k(n)=n-2$.
A two-valued polynomially witnessed cover of all these endpoints implies
$\PP^O\ne\NP^O$.
\end{corollary}
\begin{corollary}[Collapsing-oracle obstruction]
For every oracle $O$ with $\PP^O=\NP^O$, no endpoint list cover above
exists. Baker--Gill--Solovay \cite{bgs1975} supplies such oracles.
\end{corollary}
This is a theorem about matched relativizations. It does not establish
nonexistence at the empty oracle. At $O=\varnothing$, existence of the
cover remains an open sufficient condition for $\PP\ne\NP$.

\begin{obligation}[Ordinary bounded-list construction]
Fix a concrete superpolynomial $t$ and construct a verifier meeting all
three conditions for ordinary $U$. Specify the witness polynomial and
prove the global list bound. A succinct computation notation, local
transition checker, or short source program alone is not this construction.
\end{obligation}
The time function is load-bearing. The independent cost-certificate
audit gives an enumeration obstruction when $t$ is sufficiently large
relative to the verifier's witness and verification bounds. No
all-superpolynomial-$t$ construction is asserted here.

\paragraph{Verification.}
The FLMV rational inequality and bounded-list counts were checked by exact
integer/rational arithmetic on small parameters as a sanity check; the
proofs above do not rely on those finite checks. An independent REDTEAM
agent reviewed the oracle-model and witness-bound guards.
No full-repository build or proof-assistant check was performed.

\bibliographystyle{plain}
\bibliography{lower_bound_refs}
\end{document}

